% Introduction (CLAIM 2024 items 3 and 4). Every claim about prior work is a row of CLAIMS_LEDGER.md or of
% LEVELCHECK_MODEL.md section 2. No number from this study's own experiments appears here. Last paragraph =
% hypothesis and purpose. Figure 1 (the phenotypes) is cited here first, so the five figures number in the order
% of the story (1 phenotypes, 2 embedding, 3 network, 4 one exposure, 5 persons needed).
% Revision 5 (2026-09-24, REVISION_LOG3.md section 8): the LevelCheck prior-art sentences of paragraph 2 STAY (they
% motivate the problem); the fourth aim of the purpose sentence (the pre-specified comparison) moved to paper 2
% (paper2_levelcheck/MOVED_FROM_PAPER1.md item 7) and is replaced by counting the transitional persons a test needs.
% Revision 2 (2026-09-24, REVISION_LOG2.md): hsiang2011 dropped (DOI not registered; strong2021 carries the sentence);
% Mody "responding", Furuya "consecutive patients" (REFS_AUDIT.md); Huang 2022 worded from its abstract (Long-Film
% images of cadavers; the "about 17 vertebrae" and "tomosynthesis" are not in the abstract); C2 spelled out at first
% use; the network named SchwingNet (TITLE.md).
% Revision 4 (2026-09-24, Greg: thoracic level confirmation; REVISION_LOG3.md section 7): Furuya's regional rates
% (thoracic 0.42 %, lumbar 0.13 %) beside the 0.12 %, and one sentence on why a thoracic level has no
% neurophysiological check (kirshblum2011: no key muscle T2-L1 in the ISNCSCI motor examination; kulik2013: thoracic
% screw stimulation recorded from the rectus abdominis, individual-segment analysis impossible). konin2010 dropped
% with its clause (a review; nakagawa2017 carries the count consequence). Word-level trims for room, no claim changed.

\section*{Introduction}

A vertebral level is not marked; it is counted: fluoroscopy shows three or four bodies, and surgeons count
up from the sacrum or down from the second cervical vertebra (C2) across exposures
\citep{strong2021}. Half of responding spine surgeons report a wrong-level operation \citep{mody2008}, a registry
of 34,544 consecutive patients recorded 0.12\,\% (thoracic 0.42\,\%, lumbar 0.13\,\%) \citep{furuya2026}, and the level first exposed is wrong in 4.3--15\,\% of lumbar
operations before a check film corrects it \citep{dablouk2019,ammerman2006}. No neurophysiological check confirms a thoracic
level: the standard motor examination tests no key muscle between the first thoracic and second lumbar segments
\citep{kirshblum2011}, and thoracic screw stimulation, recorded from the rectus abdominis, cannot separate
segments \citep{kulik2013}. A transitional vertebra --- a
sixth lumbar (L6) or thirteenth thoracic (T13) vertebra, or a fifth lumbar vertebra (L5) fused to the sacrum
(Figure~\ref{fig:phenotypes}) --- changes the count: 8.8\,\% of 3,147 full-length radiographs lacked 24 presacral segments
\citep{sabri2025}, and in 550 lumbar operations the count was wrong in 38 of 71 patients with a lumbosacral
transitional vertebra \citep{nakagawa2017}.

Intraoperative level identification rests on registration: LevelCheck projects a reader's preoperative CT labels
onto the radiograph \citep{levelcheck2012,levelcheck2016}. Its published accuracies are
scored against those labels: 100\,\% label placement \citep{desilva2016}, and 97.6\,\%, five surgeons' median on 62
difficult films with the overlay \citep{manbachi2018}. On a transitional spine the imported
label is only as right as the planning read, and lumbar MRI localisers read without a count from C2 identified the segments in 77.9--82.2\,\% \citep{peh1999}. Registration-free naming on real C-arm images
imported identity from a preoperative radiograph (73.3\,\%) \citep{zhang2023vdvm} or needed long-length (Long-Film)
images of cadavers (88.3\,\%) \citep{huang2022}.

We hypothesised that each transitional vertebra carries local morphology that a network reading a few bodies
can learn, so a level can be named from one C-arm exposure without counting from C2 or importing a label.
The purpose was to build that network (\netname{}), examine what it learns, test it on an unseen C-arm, and
count the transitional persons a test of it needs and the public sources hold.
